\section{The family \texorpdfstring{$pq r^k$}{pq r to the k}}

\begin{theorem}
\label{thm:pqrk}
For distinct primes $p,q,r$ and every $k\geq1$, the graph $G_{pq r^k}$
has the irrational eigenvalues
\[
\frac{6k+1\pm\sqrt{(2k+2)^2-3}}2.
\]
In particular, every permutation of the exponent vector $(1,1,k)$
gives a nonintegral Laplacian spectrum.
\end{theorem}
\begin{proof}
The support classes $1,2$ have size $1$, and $13,23$ have size $k$.
Consider functions with values $s,-s$ on the first two classes,
$t,-t$ on the latter two, and zero elsewhere.
A vertex of support $1$ has $2k$ neighbors and Laplacian value
$2ks-kt$. A vertex of support $13$ has degree $4k$; after subtracting
the same-class and opposite-class values its Laplacian value is
$-s+(4k+1)t$. The swapped classes give their negatives, and all other
classes give zero by cancellation. The restricted matrix of the full
graph is therefore
\[
B_k=\begin{pmatrix}2k&-k\\-1&4k+1\end{pmatrix}.
\]
Its characteristic polynomial is $x^2-(6k+1)x+8k^2+k$, with
discriminant $(2k+2)^2-3$. Since
\[
(2k+1)^2<(2k+2)^2-3<(2k+2)^2
\]
for $k\geq1$, the two roots are irrational. Relabeling prime coordinates
preserves the graph.
\end{proof}
